\chapter*{Methven}\stepcounter{chapter}\phantomsection\addcontentsline{toc}{chapter}{Methven}

\section*{Ochil Hills}\phantomsection\addcontentsline{toc}{section}{Ochil Hills}

\section*{Seeking Aid}\phantomsection\addcontentsline{toc}{section}{Seeking Aid}
\subsection*{The Emberbound Curse}\phantomsection\addcontentsline{toc}{subsection}{The Emberbound Curse}\label{subsec:MethvenEmberboundCurse}
{\entryfont The alchemists of Methven are intrigued by the \hyperref[mi:EmberboundFragment]{\LinkFont{Emberbound Fragment}}, but their expertise is matched by their pride. They refuse to devote their time to studying the fragment without receiving something of intellectual value in return. The characters must first offer them some knowledge they consider genuinely novel - an unusual magical principle, obscure research, or insight gained through practical combat experience.

Have the contributing character explain what they offer and make an appropriate \textbf{Ability Check} based on their approach. Intelligence (Arcana) may represent magical theory, Intelligence (History) obscure knowledge, while other skills or abilities can be appropriate for knowledge gained through a character's profession or experiences.

The result determines what awaits the characters when they return several days later:}
\begin{DndTable}[header=Teach the Methven Dwarves]{lX}
	Difficulty Class	& Result \\
	13 or lower			& The dwarves dismiss the information and refuse to study the fragment. The character may try again the following day.\\
	14-16				& Their research yields 1 \hyperref[mi:PotionMaximumPower]{\LinkFont{Potion of Maximum Power}}.\\
	17+					& Two Master Alchemists take a particular interest in the fragment and each develops an experimental treatment.\\
	Natural 20			& As 17+, and the characters receive 2 Potions of Maximum Power.
\end{DndTable}
\subsubsection*{Rival Solutions}
{\entryfont When the characters return, Master Alchemist Aldr and Master Alchemist Frokia are waiting with two very different solutions. Aldr presents a fiery red ointment, while Frokia offers an icy blue preparation. Neither treatment has been tested, and neither alchemist can say with certainty what effect it will have on the fragment. Nevertheless, each loudly dismisses the other's work as amateurish and insists that only their own understanding of elemental alchemy could possibly have produced a viable solution.

Only one treatment may be chosen. Once the characters make their decision, the rejected alchemist angrily destroys their own ointment rather than allow their rival - or the characters - to benefit from it.}
\begingroup
	\DndSetThemeColor[PhbTan]
	\begin{DndComment}{DM Annotation - Experimental Ointments}
		Applying Aldr's Ointment of Fire causes the Emberbound Fragment's curse to advance by one stage. Applying Frokia's Ointment of Ice instead causes it to recede by one stage. Neither alchemist knows the true effect of their treatment before it is applied.
	\end{DndComment}
\endgroup